\usepackage[ddmmyyyy]{datetime} 
\usepackage[margin=1in]{geometry}
\usepackage{amsmath,amssymb,amsfonts} % Gộp chung các package toán
\usepackage{fancyhdr}
\usepackage{graphicx}
\usepackage{float}
\usepackage{cancel}
\usepackage{subcaption}
\usepackage[utf8]{vietnam}
\usepackage{makecell}
\usepackage{blindtext}
\usepackage{xcolor}
\usepackage{enumitem}
\usepackage{listings}
\usepackage{algorithm}
\usepackage{algpseudocode}
\usepackage{longfbox} 
\usepackage{hyperref}
\usepackage{booktabs}
\usepackage{tabularx}
\usepackage{tcolorbox}
\usepackage{colortbl}
\usepackage{multicol}
\usepackage{lipsum}

% Fix lỗi biblatex: Chỉ khai báo 1 lần duy nhất ở đây
% Nếu bạn muốn dùng style khác (như apa, ieee), hãy sửa ở phần style
\usepackage[style=numeric,backend=biber]{biblatex}

\setlength{\headheight}{30pt}

% Định nghĩa màu sắc cho Code
\definecolor{codegreen}{rgb}{0,0.6,0}
\definecolor{codegray}{rgb}{0.5,0.5,0.5}
\definecolor{codepurple}{rgb}{0.58,0,0.82}
\definecolor{backcolour}{rgb}{0.95,0.95,0.92}
\definecolor{cpptype}{rgb}{0.16, 0.47, 0.68}     % Màu xanh cho kiểu dữ liệu
\definecolor{cppfunc}{rgb}{0.6, 0.3, 0.1}        % Màu nâu/cam cho tên hàm
\definecolor{cppvar}{rgb}{0.1, 0.5, 0.5}         % Màu teal cho tên biến


% Định nghĩa Style chung
\lstdefinestyle{mystyle}{
    backgroundcolor=\color{backcolour},   
    commentstyle=\color{codegreen},
    keywordstyle=\color{magenta},
    numberstyle=\tiny\color{codegray},
    stringstyle=\color{codepurple},
    basicstyle=\ttfamily\footnotesize,
    breakatwhitespace=false,         
    breaklines=true,                 
    captionpos=b,                    
    keepspaces=true,                 
    numbers=left,                    
    numbersep=5pt,                  
    showspaces=false,                
    showstringspaces=false,
    showtabs=false,                  
    tabsize=2,
    inputencoding=utf8, % Đổi utf8x thành utf8 để ổn định hơn
    extendedchars=true,
    % Thêm ký tự đặc biệt ± để tránh lỗi compile
    literate=%
    {±}{{$\pm$}}1
    % Vần tiếng Việt (Giữ nguyên phần literate của bạn ở đây)
    {á}{{\'a}}1 {à}{{\`a}}1 {ạ}{{\d a}}1 {ả}{{\h a}}1 {ã}{{\~ a}}1
    {Á}{{\'A}}1 {À}{{\`A}}1 {Ạ}{{\d A}}1 {Ả}{{\h A}}1 {Ã}{{\~ A}}1
    {ă}{{\u a}}1 {ắ}{{\'\abreve }}1 {ằ}{{\`\abreve }}1 {ặ}{{\d \abreve }}1 {ẳ}{{\h \abreve }}1 {ẵ}{{\~\abreve }}1
    {Ă}{{\u A}}1 {Ắ}{{\'\ABREVE }}1 {Ằ}{{\`\ABREVE }}1 {Ặ}{{\d \ABREVE }}1 {Ẳ}{{\h \ABREVE }}1 {Ẵ}{{\~\ABREVE }}1
    {â}{{\^ a}}1 {ấ}{{\'\acircumflex }}1 {ầ}{{\`\acircumflex }}1 {ậ}{{\d \acircumflex }}1 {ẩ}{{\h \acircumflex }}1 {ẫ}{{\~\acircumflex }}1
    {Â}{{\^ A}}1 {Ấ}{{\'\ACIRCUMFLEX }}1 {Ầ}{{\`\ACIRCUMFLEX }}1 {Ậ}{{\d \ACIRCUMFLEX }}1 {Ẩ}{{\h \ACIRCUMFLEX }}1 {Ẫ}{{\~\ACIRCUMFLEX }}1
    {đ}{{\dj }}1 {Đ}{{\DJ }}1
    {é}{{\'e}}1 {è}{{\`e}}1 {ẹ}{{\d e}}1 {ẻ}{{\h e}}1 {ẽ}{{\~ e}}1
    {É}{{\'E}}1 {È}{{\`E}}1 {Ẹ}{{\d E}}1 {Ẻ}{{\h E}}1 {Ẽ}{{\~ E}}1
    {ê}{{\^e}}1 {ế}{{\'\ecircumflex }}1 {ề}{{\`\ecircumflex }}1 {ệ}{{\d \ecircumflex }}1 {ể}{{\h \ecircumflex }}1 {ễ}{{\~\ecircumflex }}1
    {Ê}{{\^E}}1 {Ế}{{\'\ECIRCUMFLEX }}1 {Ề}{{\`\ECIRCUMFLEX }}1 {Ệ}{{\d \ECIRCUMFLEX }}1 {Ể}{{\h \ECIRCUMFLEX }}1 {Ễ}{{\~\ECIRCUMFLEX }}1
    {í}{{\'i}}1 {ì}{{\`\i }}1 {ị}{{\d i}}1 {ỉ}{{\h i}}1 {ĩ}{{\~\i }}1
    {Í}{{\'I}}1 {Ì}{{\`I}}1 {Ị}{{\d I}}1 {Ỉ}{{\h I}}1 {Ĩ}{{\~I}}1
    {ó}{{\'o}}1 {ò}{{\`o}}1 {ọ}{{\d o}}1 {ỏ}{{\h o}}1 {õ}{{\~o}}1
    {Ó}{{\'O}}1 {Ò}{{\`O}}1 {Ọ}{{\d O}}1 {Ỏ}{{\h O}}1 {Õ}{{\~O}}1
    {ô}{{\^o}}1 {ố}{{\'\ocircumflex }}1 {ồ}{{\`\ocircumflex }}1 {ộ}{{\d \ocircumflex }}1 {ổ}{{\h \ocircumflex }}1 {ỗ}{{\~\ocircumflex }}1
    {Ô}{{\^O}}1 {Ố}{{\'\OCIRCUMFLEX }}1 {Ồ}{{\`\OCIRCUMFLEX }}1 {Ộ}{{\d \OCIRCUMFLEX }}1 {Ổ}{{\h \OCIRCUMFLEX }}1 {Ỗ}{{\~\OCIRCUMFLEX }}1
    {ơ}{{\ohorn }}1 {ớ}{{\'\ohorn }}1 {ờ}{{\`\ohorn }}1 {ợ}{{\d \ohorn }}1 {ở}{{\h \ohorn }}1 {ỡ}{{\~\ohorn }}1
    {Ơ}{{\OHORN }}1 {Ớ}{{\'\OHORN }}1 {Ờ}{{\`\OHORN }}1 {Ợ}{{\d \OHORN }}1 {Ở}{{\h \OHORN }}1 {Ỡ}{{\~\OHORN }}1
    {ú}{{\'u}}1 {ù}{{\`u}}1 {ụ}{{\d u}}1 {ủ}{{\h u}}1 {ũ}{{\~u}}1
    {Ú}{{\'U}}1 {Ù}{{\`U}}1 {Ụ}{{\d U}}1 {Ủ}{{\h U}}1 {Ũ}{{\~U}}1
    {ư}{{\uhorn }}1 {ứ}{{\'\uhorn }}1 {ừ}{{\`\uhorn }}1 {ự}{{\d \uhorn }}1 {ử}{{\h \uhorn }}1 {ữ}{{\~\uhorn }}1
    {Ư}{{\UHORN }}1 {Ứ}{{\'\UHORN }}1 {Ừ}{{\`\UHORN }}1 {Ự}{{\d \UHORN }}1 {Ử}{{\h \UHORN }}1 {Ữ}{{\~\UHORN }}1
    {ý}{{\'y}}1 {ỳ}{{\`y}}1 {ỵ}{{\d y}}1 {ỷ}{{\h y}}1 {ỹ}{{\~y}}1
    {Ý}{{\'Y}}1 {Ỳ}{{\`Y}}1 {Ỵ}{{\d Y}}1 {Ỷ}{{\h Y}}1 {Ỹ}{{\~Y}}1
}

% 1. Cấu hình mặc định cho Python (như cũ)
\lstset{style=mystyle, language=Python}

% 2. Cấu hình riêng cho C++
% 2. Cấu hình riêng cho C++ nâng cao
\lstdefinestyle{cppstyle}{
    style=mystyle,
    language=C++,
    % Nhóm 1: Từ khóa điều khiển (if, else, return...) -> Màu magenta (mặc định của bạn)
    keywordstyle=\color{magenta}\bfseries,
    % Nhóm 2: Các kiểu dữ liệu cơ bản -> Màu cpptype
    classoffset=1,
    morekeywords={int, double, float, char, bool, void, string, vector, size_t, long, long long},
    keywordstyle=\color{cpptype},
    % Nhóm 3: Các hàm phổ biến hoặc hàm tự định nghĩa -> Màu cppfunc
    classoffset=2,
    morekeywords={main, cout, cin, endl, stod, tolower, getline, push_back, substr, find, isspace, isMissingToken, trim, parseRow, computeRank},
    keywordstyle=\color{cppfunc},
    classoffset=0, % Reset lại bậc offset để không ảnh hưởng style khác
}



%%%%%%%%%%%%%%%%%%%%%%
% Cấu hình Header/Footer
\pagestyle{fancy}
\fancyhead[L]{Record Lecture Notes (AIO)}
\fancyhead[C]{}
\fancyhead[R]{\today}
\fancyfoot[L]{VNU HCMUS}
\fancyfoot[C]{}
\fancyfoot[R]{https://github.com/ThanhDanh1510}
\renewcommand{\headrulewidth}{0.4pt}
\renewcommand{\footrulewidth}{0.4pt}
%%%%%%%%%%%%%%%%%%%%%%